\documentclass[12pt]{article}
\usepackage[utf8]{inputenc}
\usepackage{listings}
\usepackage{color}
\usepackage{geometry}
\usepackage{hyperref}
\usepackage{titlesec}
\geometry{margin=1in}
\titleformat{\section}{\large\bfseries}{\thesection}{1em}{}

\definecolor{gray}{rgb}{0.5,0.5,0.5}
\definecolor{lightgray}{rgb}{0.95,0.95,0.95}

\lstset{
  backgroundcolor=\color{lightgray},
  basicstyle=\ttfamily\small,
  breaklines=true,
  frame=single,
  keywordstyle=\color{blue},
  commentstyle=\color{gray},
  showstringspaces=false
}

\title{Ejemplos de Uso del Comando \texttt{time}}
\author{}
\date{}

\begin{document}

\maketitle

El comando \texttt{time} (como \texttt{./time}, o simplemente \texttt{time}) se utiliza en sistemas Unix/Linux para medir el tiempo que tarda en ejecutarse un programa o comando. Muestra tres métricas clave:

\section*{Explicación de los tiempos obtenidos}

Al ejecutar un comando con \texttt{time}, se obtienen tres tipos de tiempo que permiten analizar el rendimiento del proceso:

\begin{itemize}
  \item \textbf{real}: también conocido como \textit{elapsed time}, es el tiempo total que transcurre desde que se inicia el comando hasta que finaliza. Incluye el tiempo de espera, ejecución, y cualquier retardo, por lo tanto puede ser afectado por otras tareas en el sistema.
  
  \item \textbf{user}: es el tiempo que la CPU dedicó a ejecutar instrucciones del programa en modo usuario (código del propio programa, sin contar llamadas al sistema).
  
  \item \textbf{sys}: es el tiempo que la CPU dedicó a ejecutar instrucciones del sistema operativo en nombre del programa, es decir, las llamadas al sistema (como acceso a archivos, entrada/salida, etc.).
\end{itemize}

Estos valores ayudan a determinar si el programa es intensivo en CPU, si hace muchas operaciones de sistema, o si pasa mucho tiempo esperando (por ejemplo, en entrada/salida o procesos externos).

\section*{Ejemplo básico}

\begin{lstlisting}
time ls -l
\end{lstlisting}

Esto mostrará cuánto tiempo tarda en ejecutarse el comando \texttt{ls -l}.

\section*{Ejemplo con un script ejecutable}

\begin{lstlisting}
chmod +x script.sh
time ./script.sh
\end{lstlisting}

Esto mide cuánto tiempo tarda en ejecutarse el script \texttt{script.sh}.

\section*{Ejemplo con un programa compilado}

\begin{lstlisting}
gcc programa.c -o programa
time ./programa
\end{lstlisting}

Esto compila y luego mide el tiempo de ejecución del binario \texttt{programa}.

\section*{Ejemplo con \texttt{sleep} (para simular una tarea)}

\begin{lstlisting}
time sleep 3
\end{lstlisting}

Resultado típico:
\begin{verbatim}
real    0m3.003s
user    0m0.000s
sys     0m0.001s
\end{verbatim}

\section*{Usando \texttt{/usr/bin/time} para más detalles}

\begin{lstlisting}
/usr/bin/time -v ./programa
\end{lstlisting}

Incluye estadísticas detalladas como uso de memoria, E/S de disco, cantidad de cambios de contexto, etc.

\end{document}
